\clearpage
\begin{proofbox}
\textbf{共同准备：先确认所有分母是否可以约掉。}

设 $p\ne0$。对 $y^2=2px$，令
\[
A=\left(\frac{a^2}{2p},a\right),\qquad
B=\left(\frac{b^2}{2p},b\right).
\]
$A,B$ 是不同的非顶点，等价于 $a\ne b$ 且 $ab\ne0$；因此 $k_{OA}=2p/a$、$k_{OB}=2p/b$ 均存在且非零。只有当 $a+b=0$ 时，$AB$ 才是竖直直线，普通斜率 $k_{AB}$ 不存在。

对 $x^2=2py$，令
\[
A=\left(a,\frac{a^2}{2p}\right),\qquad
B=\left(b,\frac{b^2}{2p}\right).
\]
同样有 $a\ne b$ 且 $ab\ne0$。这时三个斜率均存在，因为 $x_B-x_A=b-a\ne0$；但 $k_{AB}=0$ 完全可能。

\textbf{证明路线：}方法一只用初中平方差；方法二用过原点的直线斜率作参数；方法三用关于 $y=x$ 的对称理解两式；方法四用高中常规的联立与韦达定理。
\end{proofbox}

\clearpage
\begin{proofbox}
\textbf{方法一：门槛最低——初中平方差公式 + 两点斜率。}
\begin{center}% p=1, a=2,b=-1. H=(x_B,y_A). Both signed displacements are negative.
\begin{tikzpicture}[font=\small,>=stealth]
\draw[gray!45,rounded corners=3pt] (0,0) rectangle (14.2,9.0);
\node at (7.1,8.45) {平方差怎样进入弦的方向：先横移，再纵移};
\begin{scope}[shift={(2.3,3.5)},scale=1.2]
\draw[->,gray] (-1.25,0)--(3.1,0) node[above] {$x$};
\draw[->,gray] (0,-1.55)--(0,3.0) node[left] {$y$};
\draw[teal!65!black,thick,samples=100,domain=-1.4:2.4,variable=\t] plot ({\t*\t/2},\t);
\draw[blue!75!black,thin] (0,0)--(2,2);
\draw[orange!85!black,thin] (0,0)--(.5,-1);
\draw[green!45!black,very thick] (2,2)--(.5,-1);
\draw[->,violet,thick] (2,2)--(.5,2);
\draw[->,orange!85!black,thick,dashed] (.5,2)--(.5,-1);
\draw[gray] (.5,1.8)--(.7,1.8)--(.7,2);
\node[violet] at (1.25,2.4) {$\Delta x$};
\node[orange!85!black] at (-.8,.6) {$\Delta y$};
\fill (0,0) circle(2pt) node[below left=7pt] {$O$};
\fill (2,2) circle(2pt) node[below right=8pt] {$A$};
\fill (.5,-1) circle(2pt) node[right=8pt] {$B$};
\fill (.5,2) circle(1.5pt) node[above=12pt] {$H$};
\end{scope}
\node[anchor=west,align=left] at (8.0,5.0)
 {图取 $p=1,\ a=2,\ b=-1$\\[12pt]
 $A(2,2),\quad B(\frac12,-1)$\\[15pt]
 $\Delta x=\dfrac{b^2-a^2}{2p}=-\dfrac32$\\[15pt]
 $\Delta y=b-a=-3$\\[15pt]
 $\dfrac{\Delta x}{\Delta y}=\dfrac{a+b}{2p}=\dfrac12$};
\node at (7.1,.6) {两个箭头分别指向左、下；横向抛物线所需的是有向比值 $\Delta x/\Delta y$。};
\end{tikzpicture}
\end{center}

\textbf{(1) 横向开口。}设 $y^2=2px$，并采用上面的 $a,b$。由两点间横、纵坐标差，
\[
\begin{aligned}
x_B-x_A
 &=\frac{b^2-a^2}{2p}
 =\frac{(b-a)(b+a)}{2p},\\
y_B-y_A&=b-a.
\end{aligned}
\]
因为 $b-a\ne0$，所以
\[
\frac{x_B-x_A}{y_B-y_A}=\frac{a+b}{2p}.
\]
另一方面，$k_{OA}=y_A/x_A=2p/a$、$k_{OB}=2p/b$，于是
\[
\frac1{k_{OA}}+\frac1{k_{OB}}
 =\frac{a}{2p}+\frac{b}{2p}
 =\frac{a+b}{2p}
 =\frac{x_B-x_A}{y_B-y_A}.
\]
\textbf{若 $a+b\ne0$，}那么 $x_B\ne x_A$，可写为普通斜率形式
\[
\boxed{\frac1{k_{OA}}+\frac1{k_{OB}}=\frac1{k_{AB}}}.
\]
\textbf{若 $a+b=0$，}则 $AB$ 竖直，只保留上述不含 $1/k_{AB}$ 的加强式；左边和右边同为 $0$。

\end{proofbox}
\clearpage
\begin{proofbox}
\textbf{方法一（续）：纵向开口的平方差。}
\begin{center}% Coordinate interchange of C097_difference_horizontal, same units.
\begin{tikzpicture}[font=\small,>=stealth]
\draw[gray!45,rounded corners=3pt] (0,0) rectangle (14.2,9.0);
\node at (7.1,8.45) {横纵角色交换后：平方差给出普通斜率};
\begin{scope}[shift={(3.2,2.8)},scale=1.2]
\draw[->,gray] (-1.6,0)--(3.0,0) node[above] {$x$};
\draw[->,gray] (0,-.65)--(0,3.5) node[left] {$y$};
\draw[teal!65!black,thick,samples=100,domain=-1.4:2.5,variable=\t] plot (\t,{\t*\t/2});
\draw[blue!75!black,thin] (0,0)--(2,2);
\draw[orange!85!black,thin] (0,0)--(-1,.5);
\draw[green!45!black,very thick] (2,2)--(-1,.5);
\draw[->,violet,thick] (2,2)--(-1,2);
\draw[->,orange!85!black,thick,dashed] (-1,2)--(-1,.5);
\draw[gray] (-1,1.8)--(-.8,1.8)--(-.8,2);
\node[violet] at (.5,2.45) {$\Delta x$};
\node[orange!85!black] at (-1.6,1.25) {$\Delta y$};
\fill (0,0) circle(2pt) node[below left=7pt] {$O$};
\fill (2,2) circle(2pt) node[below right=8pt] {$A$};
\fill (-1,.5) circle(2pt);
\draw[gray!60] (-1.3,-.2)--(-1.03,.43);
\node at (-1.6,-.45) {$B$};
\fill (-1,2) circle(1.5pt) node[above left=7pt] {$H$};
\end{scope}
\node[anchor=west,align=left] at (8.0,5.0)
 {图取 $p=1,\ a=2,\ b=-1$\\[12pt]
 $A(2,2),\quad B(-1,\frac12)$\\[15pt]
 $\Delta x=b-a=-3$\\[15pt]
 $\Delta y=\dfrac{b^2-a^2}{2p}=-\dfrac32$\\[15pt]
 $k_{AB}=\dfrac{\Delta y}{\Delta x}=\dfrac{a+b}{2p}=\dfrac12$};
\node at (7.1,.6) {与横向图使用相同单位长度；平方差所在的位置从横移换到了纵移。};
\end{tikzpicture}
\end{center}
\textbf{(2) 纵向开口。}设 $x^2=2py$，交换横、纵坐标的角色：
\[
\begin{aligned}
k_{AB}
&=\frac{y_B-y_A}{x_B-x_A}
 =\frac{(b^2-a^2)/(2p)}{b-a}\\
&=\frac{a+b}{2p}
 =\underbrace{\frac a{2p}}_{k_{OA}}
  +\underbrace{\frac b{2p}}_{k_{OB}}.
\end{aligned}
\]
故
\[
\boxed{k_{OA}+k_{OB}=k_{AB}}.
\]
\textbf{用到的知识：}代入抛物线方程、平方差公式、直线两点斜率，初中代数即可理解。
\end{proofbox}

\clearpage
\begin{proofbox}
\textbf{方法二：步骤最少——先设 $OA,OB$ 的斜率。}
\begin{center}% y^2=2x, m=1, n=-2; second intersections A=(2,2), B=(1/2,-1).
\begin{tikzpicture}[font=\small,>=stealth]
\draw[gray!45,rounded corners=3pt] (0,0) rectangle (14.2,9.0);
\node at (7.1,8.45) {先给两条顶点连线的斜率，再确定曲线上的另一个交点};
\begin{scope}[shift={(2.1,3.5)},scale=1.25]
\draw[->,gray] (-.65,0)--(3.1,0) node[above] {$x$};
\draw[->,gray] (0,-2.2)--(0,3.1) node[left] {$y$};
\draw[teal!65!black,thick,samples=100,domain=-1.9:2.4,variable=\t] plot ({\t*\t/2},\t);
\draw[blue!75!black,very thick] (0,0)--(2.5,2.5);
\draw[orange!85!black,very thick] (0,0)--(1,-2);
\draw[green!45!black,thick] (2,2)--(.5,-1);
\node[blue!75!black] at (1.4,2.65) {$y=mx$};
\node[orange!85!black] at (1.8,-.55) {$y=nx$};
\fill (0,0) circle(2.5pt) node[below left=8pt] {$O$};
\fill (2,2) circle(2pt) node[below right=8pt] {$A$};
\fill (.5,-1) circle(2pt) node[right=8pt] {$B$};
\end{scope}
\node[anchor=west,align=left] at (8.1,5.0)
 {图取 $p=1,\ m=1,\ n=-2$\\[14pt]
 $\textcolor{blue!75!black}{A=(2p/m^2,\ 2p/m)}$\\[15pt]
 $\textcolor{orange!85!black}{B=(2p/n^2,\ 2p/n)}$\\[16pt]
 两条直线都经过 $O$\\[12pt]
 $A,B$ 各取另一个交点};
\node at (7.1,.6) {两条给定斜率的直线确定了 $A,B$，绿色弦的方向也随之确定。};
\end{tikzpicture}
\end{center}

\textbf{(1) 横向开口。}设 $k_{OA}=m$、$k_{OB}=n$，则 $m,n\ne0$。直线 $OA$ 是 $y=mx$，与 $y^2=2px$ 联立：
\[
m^2x^2=2px\quad\Longrightarrow\quad
O=(0,0),\qquad A=\left(\frac{2p}{m^2},\frac{2p}{m}\right).
\]
同理 $B=(2p/n^2,2p/n)$。当 $AB$ 非竖直时，$m+n\ne0$，从而
\[
\begin{aligned}
k_{AB}
&=\frac{2p/n-2p/m}{2p/n^2-2p/m^2}
 =\frac{(m-n)/(mn)}{(m-n)(m+n)/(m^2n^2)}\\
&=\frac{mn}{m+n}.
\end{aligned}
\]
这里 $m\ne n$，因为 $A\ne B$。于是仅需取倒数：
\[
\boxed{\frac1{k_{AB}}=\frac{m+n}{mn}=\frac1m+\frac1n}.
\]

\end{proofbox}
\clearpage
\begin{proofbox}
\textbf{方法二（续）：纵向开口用斜率确定交点。}
\begin{center}% x^2=2y, m=1,n=-1/2; A=(2,2), B=(-1,1/2).
\begin{tikzpicture}[font=\small,>=stealth]
\draw[gray!45,rounded corners=3pt] (0,0) rectangle (14.2,9.0);
\node at (7.1,8.45) {纵向模型：直线斜率直接决定另一个交点的横坐标};
\begin{scope}[shift={(3.1,2.7)},scale=1.25]
\draw[->,gray] (-1.8,0)--(3.1,0) node[above] {$x$};
\draw[->,gray] (0,-.8)--(0,3.2) node[left] {$y$};
\draw[teal!65!black,thick,samples=100,domain=-1.6:2.4,variable=\t] plot (\t,{\t*\t/2});
\draw[blue!75!black,very thick] (0,0)--(2.5,2.5);
\draw[orange!85!black,very thick] (-1.6,.8)--(0,0);
\draw[green!45!black,thick] (-1,.5)--(2,2);
\node[blue!75!black] at (1.4,2.65) {$y=mx$};
\node[orange!85!black] at (-1.0,1.7) {$y=nx$};
\fill (0,0) circle(2.5pt) node[below left=8pt] {$O$};
\fill (2,2) circle(2pt) node[below right=8pt] {$A$};
\fill (-1,.5) circle(2pt) node[above=14pt] {$B$};
\end{scope}
\node[anchor=west,align=left] at (8.2,5.0)
 {图取 $p=1,\ m=1,\ n=-\frac12$\\[15pt]
 $\textcolor{blue!75!black}{A=(2pm,\ 2pm^2)}$\\[15pt]
 $\textcolor{orange!85!black}{B=(2pn,\ 2pn^2)}$\\[16pt]
 $k_{AB}=m+n=\dfrac12$};
\node at (7.1,.6) {两条顶点连线是蓝、橙直线；连接另外两个交点得到绿色弦 $AB$。};
\end{tikzpicture}
\end{center}
\textbf{(2) 纵向开口。}同样设 $k_{OA}=m$、$k_{OB}=n$。因 $OA:y=mx$，代入 $x^2=2py$ 得另一个交点
\[
A=(2pm,2pm^2),\qquad B=(2pn,2pn^2).
\]
因此
\[
k_{AB}=\frac{2p(n^2-m^2)}{2p(n-m)}=m+n,
\]
即
\[
\boxed{k_{AB}=k_{OA}+k_{OB}}.
\]
\textbf{优点：}已知 $OA,OB$ 的斜率时，可以直接得到 $AB$ 斜率，而不必先求任意点参数；这也是处理“已知两条射线斜率，求弦斜率”题型最快的计算路线。
\end{proofbox}

\clearpage
\begin{proofbox}
\textbf{方法三：数形结合——关于 $y=x$ 对称，一次理解两个结论。}
\begin{center}% p=1, A=(1/2,1), B=(9/8,-3/2), exact reflection in y=x.
\begin{tikzpicture}[font=\small,>=stealth]
\draw[gray!45,rounded corners=3pt] (0,0) rectangle (14.2,12.0);
\node at (7.1,11.45) {关于 $y=x$ 翻转：每个点的横、纵坐标交换};
\begin{scope}[shift={(7.1,6.0)},scale=1.65]
\draw[->,gray] (-2.0,0)--(2.4,0) node[above] {$x$};
\draw[->,gray] (0,-2.0)--(0,2.4) node[left] {$y$};
\draw[gray!70,dashed] (-1.85,-1.85)--(2.25,2.25);
\draw[blue!65!black,thick,samples=110,domain=-1.85:2.05,variable=\t] plot ({\t*\t/2},\t);
\draw[orange!85!black,thick,samples=110,domain=-1.85:2.05,variable=\t] plot (\t,{\t*\t/2});
\draw[blue!75!black,very thick] (0,0)--(.5,1)--(1.125,-1.5)--cycle;
\draw[orange!85!black,very thick,dashed] (0,0)--(1,.5)--(-1.5,1.125)--cycle;
\draw[violet!55,densely dotted] (.5,1)--(1,.5) (1.125,-1.5)--(-1.5,1.125);
\fill (.5,1) circle(1.8pt) node[above left=9pt,blue!75!black] {$A$};
\fill (1.125,-1.5) circle(1.8pt) node[below=14pt,blue!75!black] {$B$};
\fill (1,.5) circle(1.8pt) node[below right=9pt,orange!85!black] {$A'$};
\fill (-1.5,1.125) circle(1.8pt) node[above left=8pt,orange!85!black] {$B'$};
\fill (0,0) circle(1.8pt);
\node at (-.65,-.3) {$O$};
\end{scope}
\node[blue!75!black] at (12.2,8.9) {$y^2=2x$};
\node[orange!85!black] at (3.1,10.3) {$x^2=2y$};
\node[gray!80!black] at (11.3,10.3) {$y=x$};
\node at (7.1,1.6) {蓝色实线三角形映成橙色虚线三角形；$O$ 保持不动。};
\node at (7.1,.65) {$\displaystyle k_{AB}=-4\quad\longleftrightarrow\quad k_{A'B'}=-\frac14$，横纵位移交换，斜率变倒数。};
\end{tikzpicture}
\end{center}

\textbf{第 1 步：镜像关系。}图中 $O$ 不动，$A,B$ 分别映到 $A',B'$。
每个点 $(x,y)$ 变为 $(y,x)$，所以 $y^2=2px$ 变为 $x^2=2py$。

\textbf{第 2 步：镜像怎样改变斜率？}对一条原斜率 $k$ 存在且不为 $0$ 的非竖直直线，其镜像的斜率是
\[
\boxed{k'=\frac{\Delta x}{\Delta y}=\frac1k}.
\]
\textbf{第 3 步：使用纵向结论。}方法一已证明镜像后的三条边满足
\[
k_{OA'}+k_{OB'}=k_{A'B'}.
\]
如果原来三条斜率都存在，这个式子就逐项变成
\[
\boxed{\frac1{k_{OA}}+\frac1{k_{OB}}=\frac1{k_{AB}}}.
\]
\textbf{读图：}对应线段的横、纵位移交换，普通斜率与倒斜率随之互换。

\textbf{边界：}原图的竖直弦在镜像后是水平弦。横向公式在该处不能写 $1/k_{AB}$，但方法一中的加强式仍成立。
\end{proofbox}

\clearpage
\begin{proofbox}
\textbf{方法四：常规解析几何——直线与抛物线联立，使用韦达定理。}
\begin{center}% y^2=2x and AB:y=2x-2, roots y=-1,2.
\begin{tikzpicture}[font=\small,>=stealth]
\draw[gray!45,rounded corners=3pt] (0,0) rectangle (14.2,10.0);
\node at (7.1,9.45) {横向开口消去 $x$：两个交点的纵坐标就是二次方程的两根};
\begin{scope}[shift={(2.3,4.0)},scale=1.25]
\draw[->,gray] (-.75,0)--(3.1,0) node[above] {$x$};
\draw[->,gray] (0,-2.25)--(0,3.0) node[left] {$y$};
\draw[teal!65!black,thick,samples=100,domain=-2.0:2.4,variable=\t] plot ({\t*\t/2},\t);
\draw[blue!75!black,thin] (0,0)--(2,2);
\draw[orange!85!black,thin] (0,0)--(.5,-1);
\draw[green!45!black,very thick] (0,-2)--(2.3,2.6);
\fill (0,0) circle(1.8pt) node[above left=7pt] {$O$};
\fill (2,2) circle(2pt) node[above left=7pt] {$A$};
\fill (.5,-1) circle(2pt) node[below=16pt] {$B$};
\end{scope}
\draw[->,gray] (9.15,1.4)--(9.15,7.7);
\node at (10.8,8.05) {$y$ 的两根};
\draw[blue!65,dashed] (4.8,6.5)--(9.15,6.5);
\draw[orange!80!black,dashed] (2.925,2.75)--(9.15,2.75);
\fill[blue!75!black] (9.15,6.5) circle(2pt);
\fill[orange!85!black] (9.15,2.75) circle(2pt);
\node[anchor=west,blue!75!black] at (9.6,6.5) {$y_A=2$};
\node[anchor=west,orange!85!black] at (9.6,2.75) {$y_B=-1$};
\node at (7.1,.65) {图取 $p=1,\ AB:y=2x-2$；联立得 $y^2-y-2=0$，根之和为 $1=2p/k$。};
\end{tikzpicture}
\end{center}

\textbf{(1) 横向开口。}设 $AB:y=kx+t$，其中 $k\ne0$ 且直线不是竖直线。把 $x=(y-t)/k$ 代入 $y^2=2px$：
\[
ky^2-2py+2pt=0.
\]
它的两个根正是 $y_A,y_B$，故由韦达定理
\[
y_A+y_B=\frac{2p}{k}.
\]
又因 $OA,OB$ 的斜率非零，
\[
\frac1{k_{OA}}+\frac1{k_{OB}}
=\frac{x_A}{y_A}+\frac{x_B}{y_B}
=\frac{y_A+y_B}{2p}
=\frac1k=\frac1{k_{AB}}.
\]
\textbf{为什么 $k\ne0$？}若横向抛物线上的两个点有相同的纵坐标，则横坐标也相同，二点重合；所以不同点不能有水平弦。

\end{proofbox}
\clearpage
\begin{proofbox}
\textbf{方法四（续）：纵向开口的两个横坐标根。}
\begin{center}% x^2=2y, AB:y=x/2+1, roots x=-1,2.
\begin{tikzpicture}[font=\small,>=stealth]
\draw[gray!45,rounded corners=3pt] (0,0) rectangle (14.2,9.0);
\node at (7.1,8.45) {纵向开口消去 $y$：改看两个交点的横坐标};
\begin{scope}[shift={(3.3,2.6)},scale=1.25]
\draw[->,gray] (-1.8,0)--(3.0,0) node[above] {$x$};
\draw[->,gray] (0,-.6)--(0,3.5) node[left] {$y$};
\draw[teal!65!black,thick,samples=100,domain=-1.6:2.5,variable=\t] plot (\t,{\t*\t/2});
\draw[blue!75!black,thin] (0,0)--(2,2);
\draw[orange!85!black,thin] (0,0)--(-1,.5);
\draw[green!45!black,very thick] (-1.6,.2)--(2.8,2.4);
\draw[blue!65,dashed] (2,2)--(2,0);
\draw[orange!80!black,dashed] (-1,.5)--(-1,0);
\fill (0,0) circle(1.8pt) node[below left=7pt] {$O$};
\fill (2,2) circle(2pt) node[above left=7pt] {$A$};
\fill (-1,.5) circle(2pt) node[above=14pt] {$B$};
\fill (2,0) circle(1.8pt) node[below=9pt,blue!75!black] {$2$};
\fill (-1,0) circle(1.8pt) node[below=9pt,orange!85!black] {$-1$};
\end{scope}
\node[anchor=west,align=left] at (8.3,4.9)
 {图取 $p=1,\ k=\dfrac12$\\[14pt]
 $AB:\ y=\dfrac12x+1$\\[15pt]
 $x^2-x-2=0$\\[15pt]
 $x_A+x_B=1=2pk$};
\node at (7.1,.6) {两个投影落在横轴的 $-1,2$ 处；根之和正是计算 $k_{OA}+k_{OB}$ 所需的量。};
\end{tikzpicture}
\end{center}
\textbf{(2) 纵向开口。}设 $AB:y=kx+t$（$k=0$ 也允许），将其代入 $x^2=2py$ 得
\[
x^2-2pkx-2pt=0.
\]
由韦达定理，$x_A+x_B=2pk$；由于 $k_{OA}=x_A/(2p)$，$k_{OB}=x_B/(2p)$，故
\[
\boxed{k_{OA}+k_{OB}=\frac{x_A+x_B}{2p}=k=k_{AB}}.
\]
\textbf{适合场景：}题目先给出弦 $AB$ 的直线方程，或需要同时研究两个交点的和、积时，韦达定理可以和原题其他条件直接衔接。
\end{proofbox}

\clearpage
\begin{proofbox}
\textbf{关联结论证明与推导一：一般二次函数的顶点平移。}
\begin{center}% Exact translation: Y=X^2 -> y=(x-2)^2+1; equal x/y units.
\begin{tikzpicture}[font=\small,>=stealth]
\draw[gray!45,rounded corners=3pt] (0,0) rectangle (14.2,10.0);
\node at (7.1,9.45) {只平移图形，三条边的方向都不变};
\node at (3.55,8.55) {$Y=X^2$};
\node at (10.65,8.55) {$y=(x-2)^2+1$};
\begin{scope}[shift={(3.2,3.8)},scale=1.15]
\draw[->,gray] (-2.1,0)--(2.1,0) node[above] {$X$};
\draw[->,gray] (0,-.75)--(0,3.4) node[left] {$Y$};
\draw[teal!65!black,thick,samples=80,domain=-1.7:1.7,variable=\t] plot (\t,{\t*\t});
\draw[blue!75!black,very thick] (0,0)--(1,1);
\draw[orange!85!black,very thick] (0,0)--(-1.5,2.25);
\draw[green!45!black,very thick] (1,1)--(-1.5,2.25);
\fill (0,0) circle(2pt) node[below left=8pt] {$V$};
\fill (1,1) circle(2pt) node[right=12pt] {$A$};
\fill (-1.5,2.25) circle(2pt) node[right=14pt] {$B$};
\end{scope}
\begin{scope}[shift={(8.1,2.65)},scale=1.15]
\draw[->,gray] (-.25,0)--(4.5,0) node[above] {$x$};
\draw[->,gray] (0,-.3)--(0,4.4) node[left] {$y$};
\draw[teal!65!black,thick,samples=80,domain=.3:3.7,variable=\t] plot (\t,{(\t-2)*(\t-2)+1});
\draw[blue!75!black,very thick] (2,1)--(3,2);
\draw[orange!85!black,very thick] (2,1)--(.5,3.25);
\draw[green!45!black,very thick] (3,2)--(.5,3.25);
\draw[gray!55,dashed] (2,0)--(2,1)--(0,1);
\fill (2,1) circle(2pt) node[below right=10pt] {$V(2,1)$};
\fill (3,2) circle(2pt) node[right=12pt] {$A$};
\fill (.5,3.25) circle(2pt) node[right=14pt] {$B$};
\end{scope}
\node at (3.55,1.8) {$k_{VA}=1,\quad k_{VB}=-\frac32$};
\node at (10.65,1.8) {$k_{VA}=1,\quad k_{VB}=-\frac32$};
\node at (7.1,.65) {两图同尺度：横、纵坐标差保持不变，故 $k_{AB}=1-\frac32=-\frac12$。};
\end{tikzpicture}
\end{center}
要判断顶点不在原点时能否使用斜率和，先配方：对 $y=ax^2+bx+c$，$a\ne0$，记
\[
h=-\frac b{2a},\qquad k=c-\frac{b^2}{4a},\qquad V=(h,k).
\]
令 $X=x-h,\ Y=y-k$，方程变成 $Y=aX^2$。坐标同时减去常数不改变两点的坐标差，因此不改变斜率。
设 $A,B$ 的新横坐标为 $X_1,X_2$，则 $X_1X_2\ne0$、$X_1\ne X_2$。于是
\[
k_{VA}=\frac{aX_1^2}{X_1}=aX_1,\qquad k_{VB}=aX_2,
\]
\[
k_{AB}=\frac{a(X_2^2-X_1^2)}{X_2-X_1}
=a(X_1+X_2)=\boxed{k_{VA}+k_{VB}}.
\]
\textbf{适用范围：}所有 $a\ne0$ 的二次函数图像都成立，向上、向下开口均可；$A,B$ 必须是不同的非顶点。
\end{proofbox}

\clearpage
\begin{proofbox}
\textbf{推导二：纵向弦斜率是两端切线斜率的平均值。}
\begin{center}% x^2=2y, p=1; A=(-1,1/2), B=(2,2); slopes -1,2,1/2.
\begin{tikzpicture}[font=\small,>=stealth]
\draw[gray!45,rounded corners=3pt] (0,0) rectangle (14.2,10.0);
\node at (7.1,9.45) {纵向：两端切线的斜率取平均，得到弦的斜率};
\begin{scope}[shift={(3.8,3.1)},scale=1.15]
\draw[->,gray] (-2.4,0)--(3.4,0) node[above] {$x$};
\draw[->,gray] (0,-1.3)--(0,4.3) node[left] {$y$};
\draw[teal!65!black,thick,samples=100,domain=-2.2:2.8,variable=\t] plot (\t,{\t*\t/2});
\draw[blue!75!black,dashed,thick] (0,0)--(-1,.5);
\draw[orange!85!black,dashed,thick] (0,0)--(2,2);
\draw[blue!75!black,very thick] (-2.2,1.7)--(.6,-1.1);
\draw[orange!85!black,very thick] (.6,-.8)--(2.8,3.6);
\draw[green!45!black,very thick] (-1,.5)--(2,2);
\fill (0,0) circle(2pt) node[below left=10pt] {$O$};
\fill (-1,.5) circle(2pt) node[above=14pt] {$A$};
\fill (2,2) circle(2pt) node[right=14pt] {$B$};
\node[blue!75!black] at (-1.6,2.1) {$t_A$};
\node[orange!85!black] at (3.25,3.25) {$t_B$};
\end{scope}
\node[anchor=west,align=left] at (8.9,5.7)
 {$x^2=2y,\quad p=1$\\[13pt]
 $A(-1,\frac12),\ B(2,2)$\\[15pt]
 $k_A=-1,\quad k_B=2$\\[15pt]
 $k_{AB}=\dfrac{-1+2}{2}=\dfrac12$};
\node at (7.1,1.3) {虚线是顶点连线，实线 $t_A,t_B$ 是切线；二者不能混淆。};
\node at (7.1,.6) {$k_{OA}=-\frac12,\ k_{OB}=1$，切线斜率分别是顶点连线斜率的 $2$ 倍。};
\end{tikzpicture}
\end{center}
本段只用联立方程的\textbf{重根}求切线，不以求导为前提。设 $x^2=2py$，$p\ne0$，端点 $A$ 的横坐标为 $u$。直线 $y=kx+d$ 与抛物线联立，得到
\[
x^2-2pkx-2pd=0.
\]
在 $A$ 处相切时，两根都等于 $u$，所以左边是 $(x-u)^2$。比较系数，得
\[
k_A=\frac up,\qquad d=-\frac{u^2}{2p}.
\]
同理，若 $B$ 的横坐标为 $v$，则 $k_B=v/p$。$A\ne B$ 保证 $u\ne v$，平方差给出
\[
k_{AB}=\frac{v^2-u^2}{2p(v-u)}
=\frac{u+v}{2p}=\boxed{\frac{k_A+k_B}{2}}.
\]
端点为顶点时对应切线斜率为 $0$，公式仍成立。
对一般二次函数，平移到顶点后，切线斜率为 $2aX_i$。由 $X_i=x_i+b/(2a)$，还原为 $2ax_i+b$；弦斜率为 $a(x_1+x_2)+b$，仍是两端切线斜率的平均值。
\end{proofbox}

\clearpage
\begin{proofbox}
\textbf{推导二（续）：横向切线取倒斜率的平均值。}
\begin{center}% Coordinate interchange of the vertical tangent diagram; identical units.
\begin{tikzpicture}[font=\small,>=stealth]
\draw[gray!45,rounded corners=3pt] (0,0) rectangle (14.2,10.0);
\node at (7.1,9.45) {横向：坐标互换后，改为两端切线的倒斜率取平均};
\begin{scope}[shift={(2.6,4.4)},scale=1.15]
\draw[->,gray] (-1.3,0)--(4.3,0) node[above] {$x$};
\draw[->,gray] (0,-2.4)--(0,3.4) node[left] {$y$};
\draw[teal!65!black,thick,samples=100,domain=-2.2:2.8,variable=\t] plot ({\t*\t/2},\t);
\draw[blue!75!black,dashed,thick] (0,0)--(.5,-1);
\draw[orange!85!black,dashed,thick] (0,0)--(2,2);
\draw[blue!75!black,very thick] (1.7,-2.2)--(-1.1,.6);
\draw[orange!85!black,very thick] (-.8,.6)--(3.6,2.8);
\draw[green!45!black,very thick] (.5,-1)--(2,2);
\fill (0,0) circle(2pt) node[below left=12pt] {$O$};
\fill (.5,-1) circle(2pt) node[right=14pt] {$A$};
\fill (2,2) circle(2pt) node[above=15pt] {$B$};
\node[blue!75!black] at (2.8,-1.7) {$t_A$};
\node[orange!85!black] at (3.6,1.95) {$t_B$};
\end{scope}
\node[anchor=west,align=left] at (8.9,5.7)
 {$y^2=2x,\quad p=1$\\[13pt]
 $A(\frac12,-1),\ B(2,2)$\\[15pt]
 $k_A=-1,\quad k_B=\frac12$\\[15pt]
 $\dfrac1{k_{AB}}=\dfrac{-1+2}{2}=\dfrac12$};
\node at (7.1,1.3) {同一对照模型交换 $x,y$；绿色弦的普通斜率变为 $k_{AB}=2$。};
\node at (7.1,.6) {这里的 $A,B$ 都不是顶点，弦不竖直，所有倒斜率才可按实数计算。};
\end{tikzpicture}
\end{center}
对 $y^2=2px$，$p\ne0$，用 $x=qy+d$ 表示切线。代入得到
\[
y^2-2pqy-2pd=0.
\]
在纵坐标为 $u$ 的 $A$ 处相切，令其等于 $(y-u)^2$，比较系数得 $q=u/p$。若 $u\ne0$，普通切线斜率 $k_A=1/q=p/u$；同理 $k_B=p/v$。

两个不同端点满足 $u\ne v$；再要求 $uv\ne0$、$u+v\ne0$，才能写
\[
\frac1{k_{AB}}=\frac{x_B-x_A}{y_B-y_A}
=\frac{u+v}{2p}
=\boxed{\frac12\left(\frac1{k_A}+\frac1{k_B}\right)}.
\]
\textbf{边界：}$u+v=0$ 时弦竖直；$u=0$ 或 $v=0$ 时有一条切线竖直。相应普通斜率不存在，原倒斜率式不能使用。
重根判定还避免了一个误解：横向抛物线与水平直线可能只有一个交点，但那条水平线并非切线。
\end{proofbox}

\clearpage
\begin{proofbox}
\textbf{推导三：抛物线焦点弦的长度倒数和。}
\begin{center}% p=2; y^2=4x. Left chord x=1; right A=(4,4),B=(1/4,-1).
\begin{tikzpicture}[font=\small,>=stealth]
\draw[gray!45,rounded corners=3pt] (0,0) rectangle (14.2,11.0);
\node at (7.1,10.45) {同一条抛物线，焦点弦方向改变，长度倒数和不变};
\node at (3.55,9.55) {竖直焦点弦};
\node at (10.65,9.55) {倾斜焦点弦};
\begin{scope}[shift={(1.6,4.9)},scale=.85]
\draw[->,gray] (-.5,0)--(5.2,0) node[above] {$x$};
\draw[->,gray] (0,-2.9)--(0,4.5) node[left] {$y$};
\draw[teal!65!black,thick,samples=100,domain=-2.7:4.25,variable=\t] plot ({\t*\t/4},\t);
\draw[blue!75!black,very thick] (1,0)--(1,2);
\draw[orange!85!black,very thick] (1,0)--(1,-2);
\fill (1,2) circle(2pt) node[right=12pt] {$A$};
\fill (1,-2) circle(2pt) node[right=12pt] {$B$};
\fill (1,0) circle(2pt) node[below right=12pt] {$F$};
\end{scope}
\begin{scope}[shift={(8.7,4.9)},scale=.85]
\draw[->,gray] (-.5,0)--(5.2,0) node[above] {$x$};
\draw[->,gray] (0,-2.9)--(0,4.5) node[left] {$y$};
\draw[teal!65!black,thick,samples=100,domain=-2.7:4.25,variable=\t] plot ({\t*\t/4},\t);
\draw[blue!75!black,very thick] (1,0)--(4,4);
\draw[orange!85!black,very thick] (1,0)--(.25,-1);
\fill (4,4) circle(2pt) node[above left=9pt] {$A$};
\fill (.25,-1) circle(2pt) node[right=14pt] {$B$};
\fill (1,0) circle(2pt) node[below right=12pt] {$F$};
\end{scope}
\node at (3.55,2.1) {$FA=FB=2$};
\node at (10.65,2.1) {$FA=5,\quad FB=\frac54$};
\node at (3.55,1.2) {$\frac12+\frac12=1$};
\node at (10.65,1.2) {$\frac15+\frac45=1$};
\node at (7.1,.6) {两图同尺度：$y^2=4x,\ p=2,\ F(1,0)$；相加的是长度倒数。};
\end{tikzpicture}
\end{center}
对 $y^2=2px$，$p>0$，焦点 $F(p/2,0)$。设过 $F$ 的直线有两个不同有限交点 $A,B$，记 $r_i=FP_i$，$P_i=(x_i,y_i)$。用距离公式，
\[
r_i^2=(x_i-p/2)^2+y_i^2=(x_i+p/2)^2.
\]
曲线上 $x_i\ge0$，所以 $r_i=x_i+p/2>0$。
焦点弦不能是 $y=0$，故可写为 $x=my+p/2$。代入得
\[
y^2-2pmy-p^2=0\quad\Longrightarrow\quad y_1y_2=-p^2.
\]
由 $x_i=y_i^2/(2p)$，得 $x_1x_2=p^2/4$。因此
\[
r_1r_2=\left(x_1+\frac p2\right)\left(x_2+\frac p2\right)
=\frac p2(x_1+x_2+p)=\frac p2(r_1+r_2).
\]
除以正数 $r_1r_2$，得 $\boxed{1/FA+1/FB=2/p}$。
竖直焦点弦是 $m=0$ 的正常情形，两段长均为 $p$，倒数和就是 $2/p$。
\end{proofbox}

\clearpage
\begin{proofbox}
\textbf{推导四：椭圆焦点弦的长度倒数和（关联阅读）。}
\begin{center}% a=5,b=3,c=4. Exact major-axis and vertical focal chords, same scale.
\begin{tikzpicture}[font=\small,>=stealth]
\draw[gray!45,rounded corners=3pt] (0,0) rectangle (14.2,9.5);
\node at (7.1,8.95) {椭圆焦点弦：两段长度改变，倒数和仍为同一个数};
\node at (3.55,8.0) {沿长轴的焦点弦};
\node at (10.65,8.0) {垂直长轴的焦点弦};
\begin{scope}[shift={(3.55,5.1)},scale=.52]
\draw[->,gray] (-5.7,0)--(5.7,0) node[above] {$x$};
\draw[->,gray] (0,-3.5)--(0,3.7) node[left] {$y$};
\draw[teal!65!black,thick] (0,0) ellipse (5 and 3);
\draw[blue!75!black,very thick] (4,0)--(5,0);
\draw[orange!85!black,very thick] (4,0)--(-5,0);
\fill (5,0) circle(2pt) node[above=20pt] {$A$};
\fill (-5,0) circle(2pt) node[above=20pt] {$B$};
\fill (4,0) circle(2pt) node[below left=10pt] {$F$};
\end{scope}
\begin{scope}[shift={(10.65,5.1)},scale=.52]
\draw[->,gray] (-5.7,0)--(5.7,0) node[above] {$x$};
\draw[->,gray] (0,-3.5)--(0,3.7) node[left] {$y$};
\draw[teal!65!black,thick] (0,0) ellipse (5 and 3);
\draw[blue!75!black,very thick] (4,0)--(4,1.8);
\draw[orange!85!black,very thick] (4,0)--(4,-1.8);
\fill (4,1.8) circle(2pt) node[above right=9pt] {$A$};
\fill (4,-1.8) circle(2pt) node[below right=9pt] {$B$};
\fill (4,0) circle(2pt) node[below left=12pt] {$F$};
\end{scope}
\node at (3.55,2.65) {$FA=1,\quad FB=9$};
\node at (10.65,2.65) {$FA=FB=\frac95$};
\node at (3.55,1.65) {$1+\frac19=\frac{10}9$};
\node at (10.65,1.65) {$\frac59+\frac59=\frac{10}9$};
\node at (7.1,.55) {两图同尺度：$\frac{x^2}{25}+\frac{y^2}{9}=1,\ F(4,0)$，定值 $\frac{2a}{b^2}=\frac{10}9$。};
\end{tikzpicture}
\end{center}
设椭圆 $x^2/a^2+y^2/b^2=1$，$a>b>0$；$a,b$ 是长、短半轴长，$c=\sqrt{a^2-b^2}$，取右焦点 $F(c,0)$。
对椭圆上的 $P(x,y)$，代入 $y^2=b^2(1-x^2/a^2)$ 得
\[
FP^2=(x-c)^2+y^2=\left(a-\frac ca x\right)^2.
\]
因为 $|x|\le a$，括号内至少为 $a-c>0$，所以 $FP=a-cx/a$。
设弦方向角为 $\theta$，焦点在弦内。令 $r=FA>0,s=FB>0$，则
\[
x_A=c+r\cos\theta,\qquad x_B=c-s\cos\theta.
\]
分别代入距离式并移项：
\[
r(a+c\cos\theta)=b^2,\qquad s(a-c\cos\theta)=b^2.
\]
两分母都不小于 $a-c>0$，所以
\[
\frac1{FA}+\frac1{FB}
=\frac{a+c\cos\theta}{b^2}+\frac{a-c\cos\theta}{b^2}
=\boxed{\frac{2a}{b^2}}.
\]
左焦点的情形由关于 $y$ 轴对称得到。这里使用的是\textbf{焦半径长度}，不是 C097 的斜率公式。
\end{proofbox}

\clearpage
\begin{proofbox}
\textbf{推导五：直角三角形的倒数平方和（初中证明）。}
\begin{center}% C=(0,0), A=(4,0), B=(0,3), D=(36/25,48/25), CD=12/5.
\begin{tikzpicture}[font=\small,>=stealth]
\draw[gray!45,rounded corners=3pt] (0,0) rectangle (14.2,9.0);
\node at (7.1,8.45) {同一个三角形，两种面积写法把斜边与高联系起来};
\begin{scope}[shift={(2.2,2.8)},scale=1.25]
\fill[blue!4] (0,0)--(4,0)--(0,3)--cycle;
\draw[blue!75!black,very thick] (0,0)--(0,3);
\draw[orange!85!black,very thick] (0,0)--(4,0);
\draw[green!45!black,very thick] (4,0)--(0,3);
\draw[violet,very thick,dashed] (0,0)--(1.44,1.92);
\draw[gray] (0,.2)--(.2,.2)--(.2,0);
\draw[gray] (1.6,1.8)--(1.48,1.64)--(1.32,1.76);
\fill (0,0) circle(2pt) node[below left=10pt] {$C$};
\fill (4,0) circle(2pt) node[below right=8pt] {$A$};
\fill (0,3) circle(2pt) node[above left=8pt] {$B$};
\fill (1.44,1.92) circle(2pt) node[above right=12pt] {$D$};
\node[blue!75!black] at (-.65,1.5) {$a=3$};
\node[orange!85!black] at (2,-.5) {$b=4$};
\node[green!45!black] at (3.15,1.25) {$c=5$};
\node[violet] at (.35,1.3) {$h$};
\end{scope}
\node[anchor=west,align=left] at (8.9,5.5)
 {$CD\perp AB$\\[14pt]
 $\frac12ab=\frac12ch$\\[14pt]
 $h=\frac{ab}{c}=\frac{12}{5}$\\[14pt]
 $\frac1{h^2}=\frac{25}{144}$};
\node at (7.1,.65) {边上标真实长度；用 $c^2=a^2+b^2$，得到 $\frac1{h^2}=\frac1{a^2}+\frac1{b^2}$。};
\end{tikzpicture}
\end{center}
设 $\angle C=90^\circ$，直角边 $BC=a>0,CA=b>0$，斜边 $AB=c$，$CD\perp AB$，$CD=h>0$。

同一个三角形的面积，既可用两条直角边计算，也可用斜边和高计算：
\[
\frac12ab=\frac12ch\quad\Longrightarrow\quad h=\frac{ab}{c}.
\]
平方并取倒数，再用勾股定理 $c^2=a^2+b^2$，得
\[
\frac1{h^2}=\frac{c^2}{a^2b^2}
=\frac{a^2+b^2}{a^2b^2}
=\boxed{\frac1{a^2}+\frac1{b^2}}.
\]
\textbf{为什么是平方的倒数？}面积关系先给出乘积比 $h=ab/c$；消去斜边时，勾股定理涉及的是平方，所以最后相加的是 $1/a^2,1/b^2$。

本段的 $a,b,c$ 是三边长，与椭圆参数无关；它是\textbf{代数结构类比}，不是抛物线结论的直接推论。
\end{proofbox}
